\uuid{Xiw6}
\titre{ Normalisation train/test : fuites de données }
\niveau{L3}
\module{Analyse de données}
\chapitre{Prétraitement des données}
\sousChapitre{Normalisation, généralisation}
\theme{StandardScaler, data leakage, fuite de données, train/test split}
\auteur{Maxime Nguyen}
\datecreate{2026-03-19}
\organisation{AMSCC}
\difficulte{3}

\contenu{
\texte{
  On compare trois façons de normaliser les données d'entraînement et de test :
  \begin{itemize}
    \item[A.] \texttt{scaler.fit\_transform(X\_train)} puis \texttt{scaler.transform(X\_test)}
    \item[B.] \texttt{scaler.fit\_transform(X\_train)} puis \texttt{scaler.fit\_transform(X\_test)}
    \item[C.] \texttt{scaler.fit\_transform(X\_all)}, puis découpage en train/test
  \end{itemize}
}
\begin{enumerate}

\item
  \question{Laquelle est correcte ?}
  \reponse{
    \textbf{A} est correcte : le \texttt{scaler} est ajusté (\texttt{fit}) uniquement sur les données
    d'entraînement, puis appliqué (\texttt{transform}) aux données de test sans les regarder.
    Cela simule fidèlement la situation réelle où les données de test sont inconnues à l'entraînement.
  }

\item
  \question{Quel problème introduit B ?}
  \reponse{
    B normalise le jeu de test avec sa propre moyenne et son propre écart-type,
    différents de ceux du jeu d'entraînement. Le modèle reçoit donc des données de test
    dans une échelle différente de celle sur laquelle il a été entraîné,
    ce qui dégrade les performances et rend l'évaluation non représentative.
  }

\item
  \question{Quel problème introduit C, plus subtil ?}
  \reponse{
    C calcule la moyenne et l'écart-type sur l'ensemble complet (train + test),
    ce qui constitue une \emph{fuite de données} (\emph{data leakage}) :
    des informations statistiques issues du jeu de test influencent la normalisation
    du jeu d'entraînement. L'évaluation est donc optimiste et ne reflète pas
    les performances réelles sur des données vraiment inédites.
  }

\end{enumerate}
}
